% !TeX root = 20260626.tex
% !LW recipe=lablog XeLaTeX
\input{_support/style.tex}

\begin{document}

\reporttitle{2026-06-26}{실험 결과 및 설계 기록}

\section*{목적과 이 보고서의 범위}

이 날짜의 실험 설정, 관찰 결과, 원인에 대한 해석과 후속 판단을 함께 정리한다. 다른 날짜의 평가 방식이나 표본 수가 다른 점수를 직접 합치지 않으며, 실행 완료와 논문 근거 채택을 구별한다.

\recordnote{수합·문서 개편: 2026-10-20. 실행 날짜는 각 절과 원시 실행 목록을 따른다. 아래의 과거 계획·진행 상태는 당시 기록이며, 현재 채택 상태와 최종 결과는 후속 실험을 함께 확인한다.}

\section{2026-06-26: OOD expert 라우팅과 문턱 조정}
\label{report:20260626_ood}
\recordnote{영어로 남아 있던 원 기록을 한국어로 수합했다. 아래 수치는 당시 설정의 기록이며 최근 TabPFN 비교와 평가 조건이 다르다.}

목적은 담당 클래스가 아닌 사례를 희소 클래스 expert가 과도하게 수락하는 현상을 줄이고,
다수 클래스 성능을 유지하면서 라우팅 구조로 소수 클래스의 탐지를 개선할 수 있는지 확인하는 것이다.
기준 실행은 \path{results/20260625_180906/}이며 원 기록의 옛 경로다.
현재 등록부에서 확인되지 않는 초기 경로는 수합 과정에서 새 실행으로 복원하거나 검증 완료로 표시하지 않았다.

첫째, expert 단위 문턱 탐색은 담당 클래스 수락과 정상·비담당 공격 수락을 함께 고려했다.
하지만 WebBot\_Rare의 잘못된 수락을 해소하지 못하고 DoS\_DDoS와 Scan\_Brute를 지나치게 보수적으로 만들어
dos-slowloris, dos-slowhttptest, ssh-patator 등의 재현율을 낮췄다.
둘째, WebBot\_Rare를 bot/web-bf/web-xss의 WebBot\_Main과
infiltration/heartbleed/web-sql의 Rare\_Micro로 분리했다.
분할 자체는 일부 라우팅을 개선했지만 작은 expert의 과도한 확신과 비담당 사례 수락은 남았다.

셋째, Rare\_Micro를 다른 expert가 모두 거절한 경우에만 호출하도록 바꿨다.
web-bf와 web-xss에 대한 잘못된 개입은 줄었지만, web-sql을 먼저 수락하는 primary expert가 있으면
Rare\_Micro에 도달할 수 없어 구제하지 못했다.
ssh-patator는 WebBot\_Main이 잘못 수락하는 문제가 남았다.
넷째, fallback margin gap을 0.5 또는 1.0으로 바꿨지만 전자는 기존 예측과 사실상 같았고,
후자는 일부 희소 클래스 선택을 전체 거절로 바꾸는 수준으로 전체 F1을 개선하지 못했다.

다섯째, validation의 정상·비담당 공격 오수락을 벌점으로 반영하는 hard-negative 문턱을 사용했다.
담당 클래스의 수락 하한은 유지했다. 전체 Macro-F1은 0.7175에서 0.7341로 변했다.
ssh-patator F1은 0.3749에서 0.5622로 개선됐지만 분리 직후 조건의 0.9841에는 못 미쳤다.
web-sql은 primary expert가 모두 거절해야만 구제 가능한 구조적 제약이 남았다.

\begin{longtable}{p{.50\linewidth}rr}
\toprule 조건 & Macro-F1 & ssh-patator F1 \\\midrule\endhead
분리 직후, fallback 제한 전 & 미기록 & 0.9841 \\
Rare\_Micro를 fallback 전용으로 제한 & 0.7175 & 0.3749 \\
fallback 전용 + hard-negative 문턱 & 0.7341 & 0.5622 \\\bottomrule
\end{longtable}

후속 개선 방향은 WebBot\_Main에 ssh-patator 같은 혼동 클래스의 음성 사례를 더 직접적으로 반영하고,
Rare\_Micro가 충분한 근거가 있을 때 primary 예측을 교정할 수 있는 제한적 경로를 검토하는 것이다.
문턱을 전체 expert에 일괄 적용하거나 fallback 선택 조건만 조정하는 것은 당시 병목을 충분히 해결하지 못했다.
이 방향은 당시 제안이며 검증된 개선 성과로 간주하지 않는다.
실행별 활성화 행렬, 클래스별 병목, 혼동행렬과 모든 실행 경로는 원 기록에 보존되어 있다.


\section*{출처와 실행 추적}

상세 수치의 원시 실행은 \path{inventory/runs.csv}, 실험별 연결은 \path{inventory/experiments.csv}에서 확인한다. 표·그림의 재현 입력은 \path{evidence/}와 \path{data/}에 둔다. 문서 이동·통합 내역은 \path{configs/experiments/migrations/20261020_documents/manifest.json}에 기록했다.

\end{document}
